\documentclass[12pt,a4paper]{article}

\usepackage[T1]{fontenc}
\usepackage[utf8]{inputenc}
\usepackage[brazil]{babel}
\usepackage{geometry}
\usepackage{booktabs}
\usepackage{tabularx}
\usepackage{longtable}
\usepackage{array}
\usepackage{hyperref}

\geometry{margin=2.5cm}
\setlength{\parskip}{0.8em}
\setlength{\parindent}{0pt}
\renewcommand{\arraystretch}{1.18}

\title{Cronograma de Execução da Pesquisa v2\\
\large Supervisor Agent for LLM Assisted Legacy Software Modernization}
\author{}
\date{Junho de 2026}

\begin{document}

\maketitle

\section{Visão Geral}

Este cronograma organiza a pesquisa em vinte e quatro meses, entre junho de 2026 e maio de 2028, com base em uma estratégia de \textit{design science research} articulada com experimentação controlada em engenharia de software \cite{hevner2004, wohlin2012, basili1986}. A proposta assume que a contribuição científica do artefato depende de três condições simultâneas: relevância do problema, qualidade do desenho do artefato e rigor na avaliação empírica. Por isso, o plano distribui o trabalho em três ciclos lógicos de pesquisa, mas duplica o esforço de coleta, análise e escrita na etapa de rigor, de modo a produzir uma primeira rodada de resultados e uma segunda rodada dedicada à replicação ou expansão do estudo.

O problema atacado é o apoio estruturado à modernização de software legado com foco em manutenibilidade, modularidade e redução de violações de projeto. Essa ênfase se alinha ao modelo de qualidade ISO/IEC 25010 e à literatura de modernização e refatoração de sistemas legados \cite{iso25010, seacord2003, feathers2004, fowler1999, martin2003}. Ao mesmo tempo, o cronograma considera o contexto recente de agentes de engenharia de software assistidos por modelos de linguagem, o que justifica a combinação entre avaliação técnica, validação humana e análise de robustez \cite{yang2024}.

Para facilitar a leitura executiva, a Tabela~\ref{tab:visao-geral} resume os grandes blocos do plano. Em seguida, cada fase é descrita em texto corrido, com ênfase nos objetivos, nos critérios decisórios e nos principais produtos esperados.

\begin{table}[htbp]
\centering
\caption{Síntese do cronograma em alto nível}
\label{tab:visao-geral}
\small
\begin{tabularx}{\textwidth}{>{\raggedright\arraybackslash}p{3.2cm} >{\raggedright\arraybackslash}p{3.2cm} >{\raggedright\arraybackslash}p{2.7cm} X}
\toprule
\textbf{Ciclo} & \textbf{Foco principal} & \textbf{Período} & \textbf{Resultado dominante} \\
\midrule
Relevância & Validação do problema e consolidação das questões de pesquisa & M1 a M4 & Questões operacionalizadas, corpus definido e protocolo pré registrado \\
Design & Endurecimento do artefato e preparação do aparato experimental & M5 a M9 & Protótipo estável, baseline funcional e piloto validado \\
Rigor I & Coleta, análise e escrita da primeira rodada empírica & M10 a M20 & Resultados confirmativos iniciais e submissão de artigo para conferência \\
Rigor II & Replicação ou expansão do estudo, consolidação final e tese & M21 a M29 & Evidência ampliada, artigo de periódico e pacote de replicação público \\
\bottomrule
\end{tabularx}
\end{table}

\section{Estrutura do Programa de Pesquisa}

O primeiro ciclo, de relevância, estabelece a fundação conceitual e metodológica da tese. Nessa etapa, o pesquisador fecha o documento de requisitos, integra a revisão sistemática à discussão de trabalhos relacionados, define um \textit{gap ledger} explícito e operacionaliza as questões de pesquisa com hipóteses, variáveis dependentes, critérios de sucesso e testes estatísticos. O encerramento dessa fase exige um corpus público defensável, um protocolo de experimento detalhado e um pré registro que congele as principais decisões analíticas antes da observação dos dados. Essa organização responde diretamente à recomendação de Hevner et al. sobre relevância do problema e rigor do processo de avaliação \cite{hevner2004}.

O segundo ciclo, de design, transforma a fundação conceitual em um artefato experimentalmente utilizável. O foco deixa de ser a formulação do problema e passa a ser a qualidade operacional do protótipo. Isso inclui controle de determinismo do cliente LLM, integração de análise estática voltada para princípios SOLID, coleta de traços completos de execução, construção da condição de baseline e implementação de critérios mínimos de testabilidade. Ao final dessa etapa, o projeto precisa demonstrar que consegue executar estudos pareados de forma reprodutível, documentada e tecnicamente auditável, o que é coerente com a tradição de experimentação em engenharia de software descrita por Wohlin et al. e Basili et al. \cite{wohlin2012, basili1986}.

O terceiro ciclo corresponde ao rigor empírico e é deliberadamente dividido em duas rodadas. A primeira rodada produz a evidência inicial necessária para um artigo de conferência, enquanto a segunda rodada reforça a validade dos achados por meio de replicação, expansão amostral ou investigação dirigida de subgrupos. Essa duplicação não é redundante. Ela reduz o risco de conclusões frágeis e fortalece tanto a contribuição científica quanto a utilidade prática do artefato.

\section{Fase 1: Relevância, definição do problema e congelamento do protocolo}

Entre junho e setembro de 2026, a pesquisa concentra-se na construção do problema científico com precisão suficiente para orientar as próximas etapas. O trabalho começa com o fechamento da especificação de requisitos em estilo Volere, seguido da integração do \textit{systematic literature review} ao capítulo de trabalhos relacionados. Em paralelo, o pesquisador documenta como o projeto atende aos sete princípios orientadores de \textit{design science research}, definindo o papel do artefato, a relevância do problema de modernização de legado, o desenho da avaliação e a forma de comunicação dos resultados \cite{hevner2004}.

No mesmo período, as questões de pesquisa são traduzidas em hipóteses observáveis e critérios de decisão. RQ1 passa a medir a capacidade do agente de produzir planos com restrições explícitas de projeto e rastreabilidade à taxonomia; RQ2 a RQ4 passam a avaliar o efeito da supervisão sobre a redução de violações SOLID em comparação com uma condição sem supervisão; RQ3 incorpora um protocolo de robustez a paráfrases; e RQ4a introduz um estudo configuracional com níveis distintos de \textit{strictness}. A definição do corpus também ocorre nessa fase, com seleção de cinco repositórios públicos em C\#, escolhidos por relevância para legado, disponibilidade de histórico e viabilidade de execução experimental.

O marco de encerramento de M4 é decisivo. A passagem para a fase seguinte só ocorre se houver questões operacionalizadas, corpus validado, desenho experimental documentado e pré registro público do protocolo. Caso contrário, a fase é estendida por algumas semanas para correção metodológica, o que é preferível a carregar ambiguidade para a etapa de implementação.

\section{Fase 2: Design, endurecimento do protótipo e piloto}

Entre outubro de 2026 e fevereiro de 2027, o cronograma desloca o foco para a consistência operacional do artefato. O primeiro bloco de trabalho trata do determinismo: inclusão explícita de \texttt{temperature = 0}, registro de \texttt{seed}, preservação de \textit{model tags} e execução de um microestudo com reexecuções controladas. O objetivo não é apenas reduzir variabilidade, mas saber medi-la, registrá-la e tratá-la como parte do próprio desenho experimental.

Em seguida, o projeto integra a camada de análise estática que permitirá mensurar mudanças associadas aos princípios SOLID. Essa camada deve ser pinada por versão, ligada a um mapeamento transparente entre regras e construtos avaliados, e acoplada a um esquema versionado de resultados por tentativa. A condição de baseline também é concluída nessa fase, preservando o mesmo modelo, o mesmo \textit{pipeline} e os mesmos mecanismos de registro, mas sem o componente de supervisão taxonômica. Essa simetria é essencial para que a comparação posterior seja interpretável.

Ainda na Fase 2, o experimento passa a registrar rastros completos de execução, incluindo repositório, revisão de código, configuração aplicada, saídas intermediárias e aprovações humanas. São acrescentados critérios de testabilidade, tais como estado de compilação, resultado dos testes e variação de cobertura, a fim de evitar leituras indevidas de ganhos em projeto quando a solução sequer mantém integridade básica de engenharia. Por fim, o piloto com dez issues mede se o protocolo de julgamento humano produz concordância aceitável, se a coleta de métricas é viável e se a etapa de paráfrases mantém equivalência semântica suficiente.

O encerramento de M9 exige seis condições: integração contínua verde, concordância inter avaliadores com $\kappa \geq 0.6$, cálculo estável de delta SOLID, reexecução controlada sem desvios não explicados, coerência do protocolo de paráfrases e cobertura mínima de design nos planos avaliados. Esse marco de decisão funciona como o verdadeiro portão de entrada para a coleta em escala.

\section{Fase 3 do Ciclo I: coleta em escala e primeira rodada analítica}

De março a agosto de 2027 ocorre a primeira coleta de dados em escala completa. O protocolo permanece congelado, e a prioridade passa a ser execução disciplinada. Para RQ1, o agente produz entre trinta e cinquenta planos, distribuídos por conjuntos de issues e acompanhados de inspeção manual sobre cobertura de SRP, OCP e DIP. Para RQ2, o estudo pareado compara supervisão e baseline em trinta a cinquenta pares, sempre com registro do estado antes e depois da intervenção. Para RQ3, o estudo de robustez expande cada tarefa em um conjunto de paráfrases semanticamente equivalentes. Para RQ4 e RQ4a, a análise se concentra respectivamente no subconjunto legado e no efeito de níveis configuracionais do supervisor.

Em setembro e outubro de 2027, os dados coletados entram primeiro em análise preliminar e depois em análise confirmativa. A análise preliminar verifica completude, valores ausentes, falhas de execução, integridade da proveniência e indícios de comportamento anômalo. Essa etapa autoriza ou não o congelamento definitivo do código de análise. A análise confirmativa, por sua vez, executa os testes definidos no protocolo, incluindo revisão de design para RQ1, Wilcoxon pareado para RQ2 e RQ4, McNemar para RQ3 e um procedimento configuracional baseado em ANOVA ou Kruskal Wallis para RQ4a, conforme a distribuição observada \cite{wohlin2012}. Em paralelo, a pesquisa realiza validação humana com pelo menos três engenheiros e conduz um instrumento de survey voltado para confiança, controle percebido e usabilidade.

O resultado esperado ao final de M19 não é apenas um conjunto de valores de $p$ e tamanhos de efeito, mas uma primeira narrativa científica coerente sobre o comportamento do artefato. Essa narrativa serve de base para a escrita do artigo de conferência, iniciada em paralelo e finalizada em M20.

\section{Fase 5 do Ciclo I: redação e submissão da primeira publicação}

Entre outubro e dezembro de 2027, a escrita do artigo da primeira rodada organiza os resultados obtidos em uma estrutura publicável. A introdução formula a relevância prática do problema de modernização assistida por modelos de linguagem; os trabalhos relacionados posicionam o artefato diante de agentes de engenharia de software, benchmarks e literatura clássica de manutenibilidade; a seção metodológica explicita o enquadramento em \textit{design science research}; e as seções de artefato, experimento e resultados consolidam o que foi efetivamente aprendido com o Ciclo I \cite{hevner2004, yang2024, iso25010}.

Essa redação tem uma função dupla. Primeiro, ela transforma resultados preliminares em contribuição científica testável. Segundo, ela expõe lacunas de explicação que precisam ser eliminadas antes da replicação. Por essa razão, o cronograma prevê revisão interna com o orientador antes da submissão, além de uma versão versionada do pacote de replicação já no fim da primeira rodada.

\section{Ciclo II: replicação, expansão da evidência e consolidação da tese}

O segundo ciclo começa ainda em novembro de 2027 com uma nova coleta, cuja estratégia é escolhida conforme os resultados do Ciclo I. A replicação pode ocorrer em um novo subconjunto de issues dos mesmos repositórios, em repositórios adicionais ou em um subgrupo particularmente relevante detectado na primeira rodada. Em qualquer uma dessas alternativas, a função do Ciclo II é separar efeito robusto de efeito circunstancial.

Entre janeiro e abril de 2028, a pesquisa repete o fluxo de limpeza, análise preliminar e análise confirmativa. O interesse aqui não está apenas em reproduzir sinais estatísticos, mas em comparar direções de efeito, magnitude das diferenças e consistência entre os dois ciclos. Sempre que possível, a síntese combina ambas as rodadas por meio de análise agregada, fortalecendo a interpretação dos achados e explicitando limites de generalização.

Entre abril e maio de 2028, a escrita final integra os dois ciclos em um artigo de periódico e em um capítulo substancial da tese. Nessa etapa, o texto deixa de relatar apenas um estudo e passa a defender uma posição científica mais madura: em que condições o supervisor melhora a qualidade do processo de modernização, como essa melhoria deve ser medida e quais ameaças à validade ainda permanecem. O pacote de replicação é então finalizado com \textit{git hashes}, versões de regras, prompts, dados anonimizados, amostra anotada e um roteiro reexecutável de ponta a ponta.

\section{Marcos decisórios e critérios de avanço}

Os principais pontos de controle do cronograma são apresentados na Tabela~\ref{tab:marcos}. Essa tabela resume a lógica de progressão do projeto: cada marco exige uma evidência mínima e explicita o que deve acontecer se tal evidência não for alcançada. Dessa forma, o plano não apenas distribui tarefas no tempo, mas também disciplina a tomada de decisão.

\begin{longtable}{>{\raggedright\arraybackslash}p{2.0cm} >{\raggedright\arraybackslash}p{4.2cm} >{\raggedright\arraybackslash}p{4.2cm} >{\raggedright\arraybackslash}p{4.2cm}}
\caption{Marcos críticos e critérios de decisão}\label{tab:marcos} \\
\toprule
\textbf{Marco} & \textbf{Critério de avanço} & \textbf{Sinal de bloqueio} & \textbf{Resposta prevista} \\
\midrule
\endfirsthead

\toprule
\textbf{Marco} & \textbf{Critério de avanço} & \textbf{Sinal de bloqueio} & \textbf{Resposta prevista} \\
\midrule
\endhead

M4 & Questões operacionalizadas, corpus validado e protocolo pré registrado & Lacunas em operacionalização ou corpus frágil & Estender a Fase 1 e corrigir o desenho \\
M9 & Piloto estável, CI verde, $\kappa \geq 0.6$ e baseline funcional & Falha de determinismo, baixa concordância ou coleta inviável & Refinar protótipo antes da coleta em escala \\
M17 & Ausência de anomalias fundamentais e código de análise pronto para congelamento & Distribuições problemáticas, dados incompletos ou validação humana insuficiente & Adiar o congelamento e documentar ajustes \\
M20 & Artigo do Ciclo I finalizado e pacote inicial de replicação preparado & Resultados pouco interpretáveis ou redação imatura & Reforçar análise e revisão antes de submeter \\
M24 & Convergência, complementaridade ou explicação clara para diferenças entre os ciclos & Contradição forte sem hipótese plausível & Investigar anomalias e recalibrar o escopo da rodada final \\
M29 & Artigo de periódico submetido, capítulo de tese consolidado e pacote público arquivado & Artefato sem estabilidade ou síntese científica incompleta & Priorizar relato transparente, inclusive de resultado negativo, antes de encerrar \\
\bottomrule
\end{longtable}

\section{Riscos metodológicos e contribuição esperada}

O cronograma responde diretamente a riscos clássicos desse tipo de pesquisa. O primeiro risco é metodológico: um enquadramento fraco de \textit{design science research} ou questões de pesquisa pouco falsificáveis podem comprometer toda a cadeia posterior. O segundo risco é instrumental: o artefato pode não ser suficientemente estável para sustentar coleta reprodutível. O terceiro risco é analítico: o estudo pode produzir efeitos pequenos, sensíveis a formulação do prompt, corpus ou configuração. A duplicação do ciclo de rigor funciona precisamente como mecanismo de mitigação desses três riscos, ao combinar pré registro, critérios explícitos de congelamento, validação humana e replicação interna \cite{hevner2004, wohlin2012, basili1986}.

Em termos de contribuição, o cronograma busca entregar mais do que um calendário operacional. Ele define uma trajetória de pesquisa em que o artefato é concebido, endurecido, avaliado e reavaliado sob critérios transparentes de qualidade. Essa estrutura torna mais defensável a tese de que um supervisor orientado por taxonomia pode apoiar a modernização de software legado de forma tecnicamente verificável e cientificamente comunicável.

\section{Conclusão}

Em síntese, o plano de vinte e quatro meses combina visão executiva e detalhamento metodológico. A tabela inicial permite enxergar o programa em alto nível, enquanto o texto corrido esclarece o papel de cada fase, os critérios de decisão e os produtos científicos esperados. O resultado é um cronograma mais adequado para uso acadêmico, para discussão com o orientador e para posterior incorporação em proposta, artigo ou capítulo de tese.

\begin{thebibliography}{99}

\bibitem{basili1986}
Basili, V. R., Selby, R. W., and Hutchens, D. H.
Experimentation in software engineering.
\textit{IEEE Transactions on Software Engineering}, 12(7):733--743, 1986.

\bibitem{feathers2004}
Feathers, M.
\textit{Working Effectively with Legacy Code}.
Prentice Hall, 2004.

\bibitem{fowler1999}
Fowler, M., Beck, K., Brant, J., Opdyke, W., and Roberts, D.
\textit{Refactoring: Improving the Design of Existing Code}.
Addison Wesley, 1999.

\bibitem{hevner2004}
Hevner, A. R., March, S. T., Park, J., and Ram, S.
Design science in information systems research.
\textit{MIS Quarterly}, 28(1):75--105, 2004.

\bibitem{iso25010}
International Organization for Standardization and International Electrotechnical Commission.
\textit{Systems and Software Engineering -- Systems and Software Quality Requirements and Evaluation (SQuaRE) -- Quality Models}.
Geneva, 2023.

\bibitem{martin2003}
Martin, R. C.
The SOLID principles.
Object Mentor, 2003.

\bibitem{seacord2003}
Seacord, R. C., Plakosh, D., and Lewis, G. A.
\textit{Modernizing Legacy Systems: Software Technologies, Engineering Processes, and Business Practices}.
Addison Wesley, 2003.

\bibitem{wohlin2012}
Wohlin, C., Runeson, P., H{"o}st, M., Ohlsson, M. C., Regnell, B., and Wessl{\'e}n, A.
\textit{Experimentation in Software Engineering}.
Springer, 2012.

\bibitem{yang2024}
Yang, J., Jimenez, C. E., Wettig, A., Lieret, K., Yao, S., Narasimhan, K., and Press, O.
SWE agent: Agent computer interfaces enable automated software engineering.
\textit{arXiv preprint arXiv:2405.15793}, 2024.

\end{thebibliography}

\end{document}